% TFM-ML-cap02-marco-contextual.tex
% Capítulo 02 — Marco Contextual
% Autor: Mario Lourido Regueira
% Scope: v5 — Corredor CHUF–Ferrol, horizonte variable
% Última revisión: 2026-04-30
% Estado: REDACTADO — §2.1, §2.2 y §2.3 alineados con v5
%         §2.4–2.5 (Caso Ferrol y Stakeholders) en espera de verificación


% ============================================================
% 2.1 EL VACÍO DE TRANSPORTE NOCTURNO EN EUROPA
% ============================================================

\section{Marco contextual y vacío de transporte nocturno}\label{cap:marco-contextual}

\subsection{El vacío de transporte nocturno en Europa}
\label{sec:vacio-nocturno}

Las ciudades europeas funcionan las veinticuatro horas del día. Los
hospitales no cierran, las cadenas de distribución no se detienen, los
hoteles no apagan sus luces. Sin embargo, los sistemas de transporte
público que conectan a sus trabajadores con sus lugares de trabajo sí lo
hacen, de forma sistemática, durante las horas de menor rentabilidad
comercial: la madrugada.

Esta discontinuidad responde a una lógica
económica que prima la eficiencia de red sobre la cobertura de necesidades:
a partir de cierta hora, la demanda cae por debajo del umbral de rentabilidad
del servicio convencional y el operador retira los vehículos. El resultado
es un vacío de entre cuatro y seis horas en el que los trabajadores de turno
nocturno deben resolver su movilidad por medios propios (vehículo
privado, taxi o, en los casos más vulnerables, a pie).

\begin{claimS}
Ese vacío se concentra típicamente entre la 01:00 y las 05:00~h
\parencite{PlyushetvaBoussauw2020}. A principios de 2019, 26 de las 28 capitales de
la Unión Europea contaban con alguna forma de transporte público nocturno,
aunque con frecuencias y coberturas muy reducidas respecto al servicio diurno
\parencite{PlyushetvaBoussauw2020}. Las ciudades no capitales (y especialmente
aquellas de menos de 150.000~habitantes) quedan en su mayor parte fuera
de este dato: sus redes nocturnas son, cuando existen, meramente testimoniales.
\end{claimS}

El vacío nocturno no es un problema homogéneo. Afecta de forma diferenciada
según tres variables: el perfil del trabajador (turno fijo o rotativo,
distancia al centro de trabajo, nivel de renta), el contexto geográfico
(ciudad grande con red nocturna deficiente frente a ciudad media sin red
alguna) y el género (la percepción de inseguridad nocturna actúa como
barrera adicional para las mujeres, que constituyen la mayoría del empleo
sanitario nocturno). El diseño del servicio propuesto en este trabajo trata
las tres variables de forma explícita.


% ============================================================
% 2.2 PANORAMA DE SOLUCIONES DE MOVILIDAD NOCTURNA
% ============================================================

\subsection{Panorama de soluciones de movilidad nocturna: modelos operativos existentes}
\label{sec:panorama-soluciones}

Antes de formular ninguna propuesta, es necesario describir
qué soluciones existen para cubrir vacíos de transporte
nocturno en corredores de demanda concentrada. Esta sección
presenta los modelos operativos documentados en la
literatura y en experiencias europeas. Ninguno se selecciona
aquí: la evaluación frente a los criterios de diseño
se realiza en el capítulo~\ref{cap:propuesta}.

\subsubsection{Transporte público convencional ampliado}
\label{sec:sol-tp-convencional}

La extensión de horarios de líneas de autobús o tren
existentes es la solución más habitual para cubrir
parcialmente la franja nocturna. Su viabilidad depende
de la existencia de demanda suficiente para justificar
el coste operativo. En ciudades de menos de 150.000
habitantes, la demanda nocturna dispersa raramente
alcanza el umbral de rentabilidad de una línea regular
\parencite{Estrada2021}.

\subsubsection{Transporte a demanda (DRT)}
\label{sec:sol-drt}

El \emph{Demand Responsive Transport} (DRT) adapta rutas
y horarios a las solicitudes de los usuarios en tiempo
real o con reserva anticipada. Reduce el coste por
kilómetro en contextos de baja densidad de demanda y
permite cubrir la franja nocturna sin la rigidez de
una línea fija. Varias ciudades europeas de tamaño medio
han implementado DRT nocturno con resultados variables
según el modelo de financiación y la tasa de adopción
\parencite{Busitalia2019QuiBus, K2Centrum2025DRT}.

\subsubsection{Shuttle con conductor en corredor fijo}
\label{sec:sol-shuttle-conductor}

Un shuttle de capacidad reducida (8--20 plazas) operando
en un corredor fijo y con horarios sincronizados con
un generador de demanda concentrada (hospital, polígono
industrial, campus logístico) es una solución de baja
complejidad tecnológica y alta adaptabilidad. No requiere
infraestructura específica ni marco regulatorio especial.
Su coste operativo está determinado principalmente por
el salario del conductor y la amortización del vehículo
\parencite{Stagecoach2020Hull}.

\subsubsection{Vehículo autónomo en corredor fijo}
\label{sec:sol-av}

Los vehículos autónomos de nivel SAE~L4 operando en
corredor definido y geofenced eliminan el coste del
conductor a largo plazo y permiten operación continua
sin límites de jornada laboral. Los pilotos europeos
más relevantes en entorno urbano (SHOW H2020, AVENUE H2020)
han demostrado viabilidad técnica en condiciones
controladas, con tasas de aceptación de usuario variables
según el contexto y la franja horaria
\parencite{SHOW2024Legacy}.
Las limitaciones actuales incluyen la madurez regulatoria,
el rendimiento en condiciones meteorológicas adversas
y el coste de la infraestructura de mapeo previo.

\subsubsection{Ride-hailing y acuerdos institucionales}
\label{sec:sol-ridehailing}

Las plataformas de \emph{ride-hailing} (VTC, taxi) y los
acuerdos de transporte institucional entre el empleador
y un operador externo representan soluciones de implantación
rápida sin necesidad de flota propia. Su coste por viaje
es generalmente superior al de un servicio de línea, pero
puede ser asumido total o parcialmente por la institución
empleadora como beneficio social \parencite{Stagecoach2020Hull, K2Centrum2025DRT}.

\subsubsection{Síntesis comparativa}
\label{sec:sol-sintesis}

Cada modelo presenta un perfil distinto de coste, madurez,
complejidad de implantación y cobertura de los criterios
de diseño que se establecerán en el capítulo~\ref{cap:criterios}.
La evaluación sistemática de estos candidatos frente a
dichos criterios se realiza en §\ref{sec:evaluacion-candidatos}.


% ============================================================
% SECCIONES PENDIENTES (cap02)
% ============================================================

% ============================================================
%  §2.3 ANÁLISIS PESTEL
%  Archivo: TFM-ML-cap02-marco-contextual.tex
%  Estado: primera redacción — pendiente revisión Mario
%  Claves .bib verificadas en sesión 2026-03-13
% ============================================================

\subsection{Análisis PESTEL: macroentorno de la movilidad nocturna en ciudades medias europeas}
\label{sec:pestel}

El análisis PESTEL permite enmarcar las condiciones de viabilidad del
servicio propuesto más allá del corredor concreto de Ferrol. Se analiza
el macroentorno de la movilidad nocturna en Europa a lo largo de seis dimensiones, con énfasis en los factores que determinan si las soluciones de transporte pueden operar de forma sostenida en ciudades de tamaño medio con vacíos de transporte nocturno documentados.

% ------------------------------------------------------------------
\subsubsection{Dimensión política}
\label{sec:pestel-politica}

La movilidad nocturna en ciudades medias europeas carece de marco 
regulatorio específico que obligue a garantizar cobertura de 
transporte público en la franja 22:00--06:00. Su provisión 
depende de la discreción del operador o la administración local, 
lo que explica la heterogeneidad de cobertura documentada entre 
capitales y ciudades de menor 
tamaño~\parencite{PlyushetvaBoussauw2020}. Esta ausencia de 
obligatoriedad es el factor político más determinante del 
macroentorno: el servicio propuesto no responde a un mandato 
legal, sino a una necesidad documentada que ningún instrumento 
vigente obliga a cubrir.

El mecanismo político-económico ordinario para sostener servicios 
de transporte que no son viables en términos comerciales es la 
Obligación de Servicio Público~(OSP), regulada en el 
Reglamento~(CE)~n.º~1370/2007. Mediante contratos OSP, la 
autoridad competente puede imponer al operador una obligación 
de cobertura e indemnizarle por el déficit generado, con el 
único requisito de que los parámetros de compensación sean 
transparentes y no produzcan sobrecompensación. Este 
instrumento es la palanca institucional primaria para que 
administraciones locales, autonómicas o el propio SERGAS puedan 
dar cobertura al servicio propuesto sin depender de la 
rentabilidad comercial de la 
franja nocturna~\parencite{Reglamento_CE_1370_2007}.

En el plano de la planificación urbana, el Reglamento~(UE) 
2024/1679 sobre la red transeuropea de transporte~(TEN-T) 
refuerza el marco político de actuación en ciudades de tamaño 
medio: designa~431 nodos urbanos europeos con obligación de 
adoptar Planes de Movilidad Urbana Sostenible~(PMUS) y remitir 
indicadores de movilidad a la Comisión~\parencite{TENT20242024}. 
Los~PMUS constituyen el instrumento de política local en el que 
puede inscribirse el servicio propuesto, articulando la cobertura 
del corredor hospitalario como medida de accesibilidad y 
bienestar laboral dentro del área funcional urbana. 
La Estrategia de Movilidad Sostenible e Inteligente de la 
Comisión proporciona el marco de financiación europeo para 
iniciativas piloto en movilidad urbana que no alcanzan 
rentabilidad comercial~\parencite{CE_Movilidad2020}.

La iniciativa URBACT \textit{Cities After Dark} visibiliza la 
movilidad nocturna como objeto de política pública municipal 
específica, y su guía de diez pasos ofrece a las ciudades de 
tamaño medio instrumentos concretos para incorporar la franja 
nocturna en sus PMUS, desde el mantenimiento de servicios 
después de las 22:00~h hasta la coordinación entre operadores 
y empleadores de turno 
nocturno~\parencite{URBACTNight2024, dAntonio2024}.

Si la tecnología seleccionada en la fase de diseño implica 
conducción automatizada, el marco político se amplía con los 
procedimientos de homologación de tipo del sistema de conducción 
automatizada establecidos por el Reglamento~(UE)~2022/1426 y 
el programa marco~ES-AV de la DGT como instrumento nacional de 
autorización de pruebas en vía 
pública~\parencite{RegUE20221426, DGTESAV2026}. Ambos 
instrumentos son condicionantes del entorno aplicables 
únicamente en ese escenario.



% ------------------------------------------------------------------
\subsubsection{Dimensión económica}
\label{sec:pestel-economica}

Los servicios de transporte nocturno en ciudades medias operan 
en un contexto económico estructuralmente adverso: la demanda 
nocturna raramente alcanza el umbral de rentabilidad comercial 
del transporte convencional, lo que obliga a retirar el 
servicio o a recurrir a financiación pública para 
sostenerlo~\parencite{Estrada2021}. El análisis de coste 
operativo en corredores de baja densidad muestra que el 
coste por pasajero de un servicio regular nocturno puede 
multiplicar entre cinco y diez veces el ingreso generado por 
tarifa, dado que los costes fijos del vehículo y del conductor 
se distribuyen sobre una demanda mínima~\parencite{UITP2024}. 
El modelo de financiación es, por tanto, una variable 
determinante con independencia de la solución técnica adoptada.

El coste laboral del conductor es la variable económica más 
determinante en cualquier solución de transporte convencional 
o semiautomatizado. En la franja nocturna, los convenios 
colectivos del sector en la mayoría de países europeos 
establecen primas sobre el salario base que incrementan 
significativamente el coste por hora de operación. Esta 
estructura de costes convierte la franja 22:00--08:00 en el 
período de mayor coste unitario del servicio y penaliza 
especialmente las soluciones con conductor 
permanente~\parencite{UITP2024}. La escasez creciente de 
conductores de transporte público en Europa~---con tasas de 
vacante que la literatura sectorial sitúa en torno al 10\% 
de los puestos y una edad media superior a los 50 
años---~añade presión adicional sobre los costes de 
contratación en horario 
nocturno~\parencite{K2Centrum2025DRT}.

El transporte a demanda~(DRT) ofrece un perfil de coste más 
favorable en contextos de baja densidad porque elimina 
kilómetros recorridos en vacío, reduce la flota necesaria 
y puede adaptar la capacidad a la demanda real de cada 
franja~\parencite{Estrada2021}. Esta ventaja es independiente 
de la tecnología de propulsión o del nivel de automatización 
del vehículo.

La demanda latente existe y tiene un coste de oportunidad 
medible. El estudio de~\textcite{PlyushetvaBoussauw2020} 
sobre trabajadores de turno nocturno en Sofía documenta que 
un empleador del sector sanitario llegaba a asumir más de 
1.020~euros mensuales en taxis nocturnos para garantizar la 
asistencia del personal. Esta cifra evidencia que la 
disposición a pagar institucional es una práctica 
documentada y que el servicio propuesto no necesita competir 
en precio con el transporte público diurno, sino con el taxi 
nocturno individual, cuya ecuación de valor es 
sustancialmente menos favorable para el empleador. El 
precedente del acuerdo entre Stagecoach y el servicio de 
salud de Hull~(Reino Unido) confirma que los modelos de 
cofinanciación empleador-administración-operador son 
viables en el contexto hospitalario 
europeo~\parencite{Stagecoach2020Hull}. La presión del 
mercado inmobiliario sobre las zonas céntricas de las 
ciudades medias desplaza progresivamente la residencia del 
personal sanitario hacia municipios periféricos, lo que 
amplía la distancia media entre domicilio y hospital y 
refuerza la demanda potencial del 
servicio~\parencite{dAntonio2024}.



% ------------------------------------------------------------------
\subsubsection{Dimensión social}
\label{sec:pestel-social}

La dimensión social del servicio está determinada por tres 
realidades estructurales que se superponen: la composición de 
género del colectivo objetivo, la percepción diferencial de 
seguridad nocturna y el riesgo documentado de fatiga al volante 
post-turno.

El empleo en el sector sanitario en la Unión Europea es 
mayoritariamente femenino: según datos de Eurostat 
correspondientes al tercer trimestre de~2021, el~78\% de los 
trabajadores del sector son 
mujeres~\parencite{Eurostat2022health}. Esta cifra es un dato 
estructural del macroentorno, no una característica local del 
caso de estudio: cualquier servicio de movilidad nocturna 
diseñado para personal sanitario tiene a las mujeres como 
usuaria principal, lo que convierte la perspectiva de género 
en un requisito de diseño de primer 
orden~\parencite{Eurostat2020Health}.

La percepción de seguridad nocturna muestra una brecha de 
género cuantificada y consistente en múltiples contextos 
europeos. \textcite{PlyushetvaBoussauw2020} midieron en Sofía 
que los hombres valoraban su seguridad nocturna en~5,82 sobre 
10 frente a~4,42 de las mujeres, diferencia que se traduce 
directamente en menor disposición a utilizar el transporte 
público durante la franja nocturna y en mayor dependencia de 
alternativas privadas de mayor coste. Un análisis de datos 
del metro y autobús de Londres entre 2009 y~2018 confirma 
que las mujeres se sienten entre un~6\% y un~10\% más 
inseguras que los hombres en el transporte público, con 
independencia del 
modo~\parencite{AitBihiOuali2020}. Esta brecha no es un 
factor de confort: condiciona la elección modal y determina 
si el servicio es o no adoptado por su usuaria principal.

El trabajo a turnos, y especialmente el turno nocturno, está 
documentado como factor de riesgo para la salud física y 
mental. La alteración del ritmo circadiano genera déficit 
de sueño crónico, mayor incidencia de fatiga, síntomas 
depresivos y deterioro cognitivo acumulado, con efectos 
que se agravan en patrones de turno rotatorio e irregular 
\parencite{SanchezSellero2025}. En el sector sanitario, 
donde la complejidad y la carga cognitiva de la jornada 
son elevadas, esta fatiga se extiende al desplazamiento 
post-turno: el trayecto de regreso a casa al amanecer 
constituye la ventana de mayor riesgo de accidente de 
tráfico por somnolencia al volante, coincidiendo con el 
punto más bajo del ciclo circadiano~\parencite{Eurofound2020Sectors}. 
Este riesgo es independiente de la tecnología de transporte: 
afecta a cualquier trabajador que conduzca su propio vehículo 
en esa franja, y su eliminación requiere una alternativa 
colectiva disponible en ese horario concreto.

El perfil demográfico del colectivo refuerza la urgencia del 
factor anterior. Más del~36\% de los trabajadores sanitarios 
de la UE tienen~50 o más años, y en quince de los 
veinticuatro países para los que existen datos, más del~30\% 
de las enfermeras en ejercicio supera los~55 
años~\parencite{Eurostat2022health}. Esta concentración de 
trabajadoras de mayor edad en el turno nocturno amplifica 
la vulnerabilidad fisiológica a la fatiga y refuerza la 
necesidad de soluciones de movilidad que no dependan de la 
capacidad de conducción individual.



% ------------------------------------------------------------------
\subsubsection{Dimensión tecnológica}
\label{sec:pestel-tecnologica}

El paisaje tecnológico disponible para servicios de transporte 
nocturno en corredores de baja demanda abarca cuatro familias 
con niveles de madurez distintos: transporte convencional en 
vehículo de capacidad reducida, plataformas de transporte a 
demanda (DRT), electrificación de flota y vehículos con 
conducción automatizada. Esta sección describe su estado como 
factores del entorno. La evaluación de cada familia frente 
a los criterios del servicio se realiza en~§\ref{sec:evaluacion-tecnologica}.

El transporte convencional en minibús o autobús de capacidad 
reducida es la familia de mayor madurez operativa. Sus 
principales innovaciones recientes se concentran en la 
integración de trenes motrices eléctricos e híbridos y en 
sistemas ITS embarcados~---información en tiempo real, 
videovigilancia, asistencia al conductor---~más que en el 
vehículo base~\parencite{UITP2024}.

El transporte a demanda ha evolucionado de sistemas de reserva 
telefónica a plataformas digitales que utilizan algoritmos de 
optimización de rutas y asignación de vehículos en tiempo real 
o con reserva anticipada. Esta madurez lo convierte en la 
opción más documentada para corredores de baja densidad fuera 
del pico: el análisis de coste operativo demuestra que el DRT 
es preferible al servicio regular cuando la demanda es dispersa, 
al reducir el número de kilómetros recorridos en 
vacío~\parencite{Estrada2021}. Una revisión de servicios DRT 
activos en Europa documenta su despliegue bajo modelos de 
financiación y cobertura muy 
heterogéneos~\parencite{K2Centrum2025DRT}.

La electrificación de la flota es una variable tecnológica 
independiente del nivel de automatización del vehículo. Puede 
aplicarse a cualquiera de las familias anteriores y constituye 
la vía de descarbonización más inmediata del servicio con 
independencia del modo de conducción 
adoptado~\parencite{UITP2024}.

Los vehículos con conducción automatizada de nivel SAE~L4 
representan la familia de menor madurez para despliegue 
comercial estable. Los pilotos europeos han demostrado 
viabilidad técnica en entornos controlados, con tasas de 
aceptación de usuario variables según el contexto y la 
franja horaria~\parencite{Debbaghi2025}. Las limitaciones 
documentadas del entorno incluyen la madurez regulatoria, 
el rendimiento en condiciones meteorológicas adversas y la 
ausencia de despliegues comerciales a gran escala en el 
horizonte inmediato~\parencite{KonstantasAVENUE2023lessons, 
DibaGore2025}. Estas limitaciones son factores del macroentorno 
que cualquier propuesta con componente de automatización debe 
incorporar como condicionantes externos, no como problemas 
a resolver internamente.

% ------------------------------------------------------------------
\subsubsection{Dimensión ecológica}
\label{sec:pestel-ecologica}

La ausencia de transporte público nocturno en ciudades medias 
europeas tiene un coste ecológico cuantificable: cada 
trabajador nocturno sin alternativa colectiva genera un 
desplazamiento individual en vehículo privado, el modo de 
mayor emisión por viajero-kilómetro. En el corredor 
CHUF--Ferrol, los datos de la encuesta primaria confirman 
que el vehículo propio es la alternativa modal predominante 
en la franja nocturna~\parencite{LouridoRegueira2026}.

La sustitución de esos desplazamientos individuales por un 
servicio colectivo encaja directamente en el marco del Pacto 
Verde Europeo, cuyo objetivo de neutralidad climática para 
2050 exige reducir las emisiones del transporte en un~90\% 
respecto a los niveles de~1990~\parencite{CE_Movilidad2020}. 
El punto de partida no es neutro: el transporte colectivo 
convencional ya es intrínsecamente más eficiente que el 
vehículo privado, y esta ventaja se acumula cuando la 
comparación se establece frente al taxi nocturno individual, 
alternativa modal predominante en ausencia de servicio 
público~\parencite{UITP2024}.

La electrificación de la flota es un factor ecológico 
independiente del nivel de automatización del vehículo. 
La Directiva~(UE)~2019/1161 sobre vehículos limpios 
establece cuotas mínimas de adquisición de vehículos de 
bajas emisiones para las flotas de transporte público: 
el~45\% de las compras de autobús deben cumplir criterios 
de limpieza en el período~2021--2025, umbral que sube 
al~65\% en el período~2026--2030~\parencite{CVD20191161}. 
Este marco regulatorio implica que cualquier solución de 
transporte colectivo que entre en servicio en el horizonte 
del proyecto deberá incorporar propulsión eléctrica o de 
bajas emisiones con independencia de si el vehículo dispone 
de conducción automatizada. La electrificación no es una 
característica tecnológica diferencial de una candidata 
concreta: es una condición del entorno regulatorio 
aplicable a todas ellas.

Para un corredor nocturno de demanda concentrada y horarios 
predecibles, la electrificación permite además planificar 
ciclos de recarga sincronizados con los períodos de baja 
demanda entre cambios de turno, optimizando el consumo sin 
requerir infraestructura de recarga 
rápida~\parencite{UITP2024}.



% ------------------------------------------------------------------
\subsubsection{Dimensión legal}
\label{sec:pestel-legal}

El marco legal del servicio propuesto opera en tres planos 
independientes de la tecnología que se seleccione: los derechos 
de los trabajadores nocturnos, los instrumentos de contratación 
pública para servicios de transporte por debajo del umbral 
comercial, y la protección de datos en sistemas de gestión 
digital de la demanda.

En el plano laboral, la Directiva~2003/88/CE sobre la ordenación 
del tiempo de trabajo establece protecciones específicas para los 
trabajadores nocturnos: límite de ocho horas medias de trabajo 
nocturno por período de~24 horas, derecho a evaluaciones de salud 
periódicas gratuitas y posibilidad de traslado a turno diurno 
cuando surjan problemas de salud vinculados al 
turno~\parencite{Directiva200388CE}. Esta protección normativa 
refuerza el argumento de que las instituciones públicas 
---SERGAS, Concello de Ferrol, Xunta de Galicia--- tienen 
fundamento legal para priorizar la movilidad del personal de 
turno nocturno dentro de sus políticas de bienestar laboral. 
La Directiva no regula el desplazamiento al trabajo, pero 
reconoce el carácter estructuralmente más oneroso de la 
jornada nocturna: ese reconocimiento es el punto de partida 
institucional del servicio propuesto.

En el plano de la contratación pública, el 
Reglamento~(CE)~n.º~1370/2007 sobre servicios públicos de 
transporte de viajeros define el mecanismo central para 
sostener servicios que no son viables en términos comerciales: 
la Obligación de Servicio Público~(OSP). Mediante contratos OSP, 
la autoridad competente puede imponer al operador una obligación 
de servicio e indemnizarle por el déficit generado, siempre que 
los parámetros de compensación sean transparentes y no 
sobrecompensen al operador~\parencite{Reglamento_CE_1370_2007}. 
Para un servicio nocturno en un corredor de baja demanda como 
el corredor~CHUF, la OSP no es una solución excepcional sino 
el instrumento ordinario: la literatura sectorial confirma que 
la viabilidad comercial autónoma de servicios nocturnos en 
ciudades de menos de~150.000 habitantes es 
estructuralmente improbable~\parencite{UITP2024}.

En el plano de la protección de datos, cualquier servicio que 
gestione la demanda de forma digital ---ya sea mediante 
reserva anticipada o asignación dinámica--- procesa datos de 
localización y patrones de desplazamiento de los usuarios. 
El Reglamento~(UE)~2016/679~(RGPD) exige base jurídica 
legítima para ese tratamiento, minimización de datos, 
limitación de finalidad y medidas técnicas y organizativas 
proporcionales al riesgo~\parencite{GDPR20162016679}. Para 
un servicio orientado a trabajadores de una institución 
pública, el interés legítimo o la ejecución de un contrato 
constituyen bases jurídicas más sólidas que el consentimiento, 
dada la relación de dependencia laboral. Este factor es 
independiente de la tecnología de transporte y aplica a 
cualquier variante del servicio que incorpore gestión digital 
de la demanda.

Si la tecnología seleccionada en la fase de diseño implica 
conducción automatizada, el marco legal se amplía con el 
régimen de responsabilidad por productos defectuosos, 
updated por la Directiva~(UE)~2024/2853, que extiende 
el concepto de producto defectuoso al software y a los 
sistemas de inteligencia artificial~\parencite{DirectivaUE20242853}, 
y con el procedimiento de homologación de tipo del sistema 
de conducción automatizada~(ADS) establecido por el 
Reglamento~(UE)~2022/1426~\parencite{RegUE20221426}. Ambos 
instrumentos son condicionantes del entorno aplicables 
únicamente en ese escenario.

\subsection{Caso de estudio: corredor CHUF--Ferrol}
\label{sec:caso-ferrol}

Ferrol es un municipio gallego de \num{64358}~habitantes
\parencite{INE2024Ferrol} situado en el extremo noroeste de la
Península Ibérica, con una estructura urbana dispersa y una red de
transporte público históricamente orientada a la franja diurna. El
Complexo Hospitalario Universitario de Ferrol (CHUF) concentra la
actividad sanitaria de referencia del Área Sanitaria de Ferrol y se
ubica en el polígono de A~Cabana, en la periferia occidental del
municipio \parencite{Sanidad2024CHUF}, a una distancia estimada de entre
4,2 y 5,2\,km del centro urbano por las vías de acceso habituales
\parencite{GoogleMaps2026Ferrol}.
El \'{A}rea Sanitaria de Ferrol cuenta con aproximadamente
2.300~profesionales \parencite{SERGASFerrol2025}, de los
cuales una parte opera en turno nocturno en el~CHUF.
La franja comprendida entre las 22:00 y las 06:00~h carece de servicio
de transporte público regular en el corredor CHUF--zonas residenciales
\parencite{Alsaferrol2024horarios,RenfeCercaniasFerrol2026}. Los cambios
de turno del personal sanitario se producen a las 22:00 y las 08:00~h
\parencite{XuntaGalicia2025dog}, coincidiendo con la franja de mayor vacío de oferta y con el periodo 
de mayor riesgo de fatiga al volante documentado en la literatura 
sobre trabajo nocturno \parencite{SanchezSellero2025, Eurofound2020Sectors}.



Ferrol actúa en este trabajo como laboratorio de diseño representativo
de un nicho sistemáticamente desatendido en Europa: ciudades de entre
50.000 y 150.000~habitantes con un hospital de referencia en
localización periférica y un vacío de transporte nocturno estructural.
El modelo de servicio generado desde este caso está explícitamente
concebido para ser transferible a corredores análogos; su capacidad de
replicabilidad se desarrolla en el capítulo~\ref{cap:replicabilidad}.


